\documentclass[11pt]{article}
\usepackage[margin=1in]{geometry}
\usepackage{amsmath,amssymb}
\usepackage[hidelinks]{hyperref}

\begin{document}

\section*{Human Review: Sharp Constant for Proposition P:schur}

\noindent Original automatic proof: \texttt{solution\_packet.pdf}

\noindent Checked date: 15 June 2026

\noindent Status: Verified.

\noindent Verdict: The proof appears correct. It gives a sharpness example showing that, for \(c\geq 1\), the constant \(2c+1\) in Kalenda's implication
\[
(iii_c)\Longrightarrow(ii_{2c+1})
\]
cannot be improved.

\section*{Human Comments}

The construction is well targeted. For \(c\geq1\), the proof considers
\[
X_c=\mathbb K\oplus\ell_1
\]
with norm
\[
\|(t,x)\|_c=\max\{\|x\|_1,\ |t|+2c\|x\|_\infty\}.
\]
The sequence
\[
y_n=(1,e_n)
\]
then has
\[
\|y_n\|_c=2c+1.
\]
On the other hand, the averaging argument shows that every weak-star cluster point of \((y_n)\) in \(X_c^{**}\) has norm exactly \(1\). Hence condition \((i_C)\) fails for every \(C<2c+1\). Since Proposition P:schur in the source paper proves \((i_C)\Longleftrightarrow(ii_C)\), this also gives failure of \((ii_C)\) for every \(C<2c+1\). This proves sharpness in the formulation of Question 9.2.

The proof that \(X_c\) satisfies \((iii_c)\) also looks correct. The key point is to take an arbitrary bounded sequence \(z_n=(t_n,x_n)\), pass to a subsequence with \(t_n\to t\), \(x_n\) converging coordinatewise to some \(x\in\ell_1\), and \(\|x_n-x\|_1\to A\). The fact that the coordinatewise limit belongs to \(\ell_1\) follows from Fatou's lemma, after the usual diagonal extraction. The gliding-hump lemma then gives
\[
\delta_{\ell_1}(x_n-x)\geq 2A
\]
and, after passing to a further subsequence, the estimates
\[
\|u_p-u_q\|_1\leq 2A+\varepsilon,
\qquad
\|u_p-u_q\|_\infty\leq A+\varepsilon,
\]
where \(u_n=x_n-x\). These estimates give
\[
\widetilde{\operatorname{ca}}(z_n)\leq 2cA
\leq c\,\delta_{X_c}(z_n),
\]
which is precisely \((iii_c)\).

The only improvement I would suggest is in the proof of Lemma 1. The argument uses the fact that \(\ell_1\) has the \(1\)-Schur property, namely
\[
\operatorname{ca}(w_n)\leq \delta_{\ell_1}(w_n)
\]
for bounded sequences in \(\ell_1\). This is a standard quantitative Schur fact, but it should ideally be cited explicitly. In the source paper, this is recalled in Example 5.9, where \(\ell_1\) is said to have the \(1\)-Schur property by Kalenda--Spurny, Theorem 1.3. Adding this citation at the use in Lemma 1 would make the proof more self-contained and easier to check.

Aside from this minor citation/reference issue, the main construction, the verification of \((iii_c)\), the cluster-point norm calculation, and the final sharpness argument all look sound.

\section*{Review Outcome}

Accepted as an essentially correct solution to the optimality question for \(c\geq1\). The small-constant range \(0<c<1\) is also plausibly handled by the Rosenthal--James argument in the packet, but the main contribution for Question 9.2 is the sharp \(2c+1\) example above.

\end{document}
